% ============================================================
% diff 报告 — 初始提交 → stab27 全量改动概览（PDF 版）
% 编译: xelatex -interaction=nonstopmode -halt-on-error -file-line-error "diff_init_stab27_20260818.tex"
% 交付前 Overfull 计数必须为 0
% ============================================================
\documentclass[UTF8]{ctexart}
\usepackage[margin=2.5cm]{geometry}
\usepackage{graphicx}
\usepackage{amsmath}
\usepackage{microtype}
\usepackage{listings}
\usepackage{xcolor}
\usepackage{booktabs}
\usepackage{longtable}
\usepackage{xltabular}
\usepackage{tabularx}
\usepackage{hyperref}
\usepackage{fancyhdr}
\usepackage[most]{tcolorbox}
\usepackage{titlesec}

\definecolor{accent}{HTML}{1F4E79}
\definecolor{evidence}{HTML}{1E8449}
\definecolor{reference}{HTML}{8C6D1F}
\definecolor{conclusion}{HTML}{8B1A1A}
\definecolor{codebg}{HTML}{F4F6F7}
\definecolor{struct}{HTML}{3B3B98}
\definecolor{relation}{HTML}{0E7C7B}
\definecolor{idea}{HTML}{B03A2E}
\definecolor{warn}{HTML}{D35400}
\definecolor{hl}{HTML}{FDF6E3}

\setlength{\emergencystretch}{3em}
\hypersetup{colorlinks=true,linkcolor=accent,urlcolor=relation}

\titleformat{\section}{\Large\bfseries\color{accent}}{\thesection}{1em}{}
\titleformat{\subsection}{\large\bfseries\color{struct}}{\thesubsection}{1em}{}
\titleformat{\subsubsection}{\normalsize\bfseries\color{struct!70!black}}{\thesubsubsection}{1em}{}

\newtcolorbox{conclbox}{breakable,colback=conclusion!6,colframe=conclusion,title=结论}
\newtcolorbox{refbox}{breakable,colback=reference!6,colframe=reference,title=参考背景}
\newtcolorbox{ideabox}{breakable,colback=idea!6,colframe=idea,title=项目思路}
\newtcolorbox{warnbox}{breakable,colback=warn!6,colframe=warn,title=局限与风险}
\newtcolorbox{keybox}{breakable,colback=hl,colframe=accent,title=要点}

\lstset{
  basicstyle=\ttfamily\footnotesize,
  backgroundcolor=\color{codebg},
  frame=single,
  firstnumber=auto,
  keywordstyle=\color{accent}\bfseries,
  commentstyle=\color{evidence!70!black},
  stringstyle=\color{struct},
  breaklines=true,
  breakatwhitespace=false,
  postbreak=\mbox{\textcolor{accent}{$\hookrightarrow$}\space},
  keepspaces=true,
  columns=flexible,
  numbers=left,numberstyle=\tiny\color{struct!60!black},
}

\title{PyQCU 改动概览（init → stab27）\\[6pt]\large 初始提交 0204fa5 → 最新标签 stab27 全量 git diff 结构化分析}
\author{opencode diff}
\date{2026-08-18}

\begin{document}
\maketitle

\begin{abstract}
本报告对 PyQCU 从初始提交 \texttt{0204fa5} 到最新标签 \texttt{stab27} 的全部改动做只读结构化概览。全量 2138 文件、+1,163,982/−2 行、区间内 1120 提交；其中外部库镜像 \texttt{refer/}（1559 文件）与测试/产物 \texttt{examples/}、\texttt{logs/} 占增行绝大部分，真正自有库代码集中在 \texttt{pyqcu/}（+6847）与 \texttt{cpp/}（+16234）。边界检查：冲突标记 0、二进制 170、未检测重命名。
\end{abstract}

\tableofcontents

\section{改动概览}

\begin{keybox}
\textbf{范围}: 初始提交 \texttt{0204fa5} (init) → 标签 \texttt{stab27}（全历史）\\
\textbf{提交}: 区间内 1120 个提交\\
\textbf{统计}: 2138 个文件, +1,163,982 / -2 行\\
\textbf{分组}: 新增(全量) / 外部镜像 \texttt{refer/} 1559 文件 / 测试与产物 \texttt{examples}+\texttt{logs} 384 文件\\
\textbf{主导}: \texttt{refer/git-rep/}（外部库镜像）与 \texttt{examples/pyqcu/*.txt}（运行日志）占增行绝大部分\\
\textbf{自有库}: \texttt{pyqcu/} 47 文件 +6847 行；\texttt{cpp/} 78 文件 +16234 行\\
\textbf{关注}: 冲突标记 0；二进制 170；未检测重命名
\end{keybox}

\section{规模与分布（实测）}

\begin{xltabular}{\textwidth}{lrrX}
\toprule
\textcolor{struct}{顶层目录} & 文件数 & 增行 & 说明 \\
\midrule
\endhead
\texttt{refer} & 1559 & 385,023 & 外部库镜像（DDalphaAMG/PyQUDA/quda 等），非自有代码 \\
\texttt{examples} & 136 & 723,046 & 测试入口与运行日志（\texttt{*.txt} 占主导） \\
\texttt{logs} & 261 & 29,394 & 测试/调试产物归档（\texttt{.gitignore} 豁免但历史已提交） \\
\texttt{cpp} & 78 & 16,234 & C++ CUDA 后端（手写内核 + MPI halo） \\
\texttt{pyqcu} & 47 & 6,847 & 纯 Python 实现 + Cython 桥 + cann 兼容层 \\
\texttt{.opencode} & 36 & 2,493 & opencode 运行环境，与库无关 \\
\texttt{docs} & 12 & 730 & 文档与既有分析报告 \\
根文件 & 8 & 215 & \texttt{setup.py}/\texttt{env.sh}/\texttt{build.sh}/\texttt{AGENTS.md} 等 \\
合计 & 2138 & 1,163,982 & — \\
\bottomrule
\caption{全量 diff 规模按顶层目录（数据来源: git diff --numstat 0204fa5 stab27 实测）}
\end{xltabular}

\section{PyQCU 自有库代码演进}

真正自有库代码集中在 \texttt{pyqcu/} 与 \texttt{cpp/}。\texttt{pyqcu/} 子目录分布：

\begin{table}[htbp]\centering
\begin{tabular}{lrr}
\toprule
子目录 & 文件数 & 增行 \\
\midrule
\texttt{tools} & 9 & 1846 \\
\texttt{cuda} & 11 & 1604 \\
\texttt{dslash} & 5 & 1099 \\
\texttt{solver} & 5 & 749 \\
\texttt{testing} & 2 & 697 \\
\texttt{smear} & 3 & 248 \\
\texttt{cann} & 4 & 286 \\
\texttt{lattice} & 2 & 250 \\
其他 & 7 & 68 \\
\bottomrule
\end{tabular}
\caption{pyqcu/ 子目录分布（数据来源: git diff --numstat 实测）}
\end{table}

\section{主导文件（自有库，按增行）}

\begin{xltabular}{\textwidth}{lXr}
\toprule
\textcolor{struct}{文件} & 角色 & 增行 \\
\midrule
\endhead
\texttt{\nolinkurl{cpp/cuda/qcu/include/lattice_clover_multigrid.h}} & Clover 多重网格核心头 & 1833 \\
\texttt{\nolinkurl{cpp/cuda/qcu/src/clover_dslash_multi.cu}} & 多重网格 Clover dslash 内核 & 1610 \\
\texttt{\nolinkurl{cpp/cuda/qcu/src/wilson_dslash.cu}} & Wilson dslash 内核 & 1442 \\
\texttt{\nolinkurl{cpp/cuda/qcu/src/multigrid.cu}} & 多重网格 V-cycle 内核 & 981 \\
\texttt{\nolinkurl{cpp/cuda/qcu/src/clover_dslash_single.cu}} & 单层 Clover dslash 内核 & 918 \\
\texttt{\nolinkurl{pyqcu/tools/_multigrid.py}} & 纯 Python 多重网格工具 & 819 \\
\texttt{\nolinkurl{pyqcu/testing/__init__.py}} & 集成测试套件 & 658 \\
\texttt{\nolinkurl{pyqcu/cuda/_multi_gpu.py}} & 一线程一卡多线程多卡 MG & 568 \\
\texttt{\nolinkurl{pyqcu/solver/_multigrid.py}} & 纯 Python 多重网格求解器 & 558 \\
\texttt{\nolinkurl{pyqcu/tools/_define.py}} & 参数协议常量定义 & 501 \\
\bottomrule
\end{xltabular}

\section{边界检查（易漏项）}

\begin{itemize}
  \item \textbf{冲突标记}：\texttt{git diff | grep '^<<<<<<<|^=======|^>>>>>>>'} → \textcolor{evidence}{0}（无合并冲突残留）。
  \item \textbf{调试残留}：主导增量文件为 \texttt{examples/pyqcu/log-v*.txt}（运行日志），无库代码内 \texttt{print/echo} 调试行确证。
  \item \textbf{未跟踪文件}：本视图为版本间 diff，不含未跟踪内容。
  \item \textbf{二进制/大文件}：\texttt{git diff --numstat} 首字段为 \texttt{-} 的共 \textcolor{evidence}{170} 个，主要为 \texttt{refer/git-rep/} 外部库二进制与 \texttt{examples/} 数据/\texttt{.ipynb}。
  \item \textbf{权限位变化}：\texttt{--summary} 未报告 mode 变化。
  \item \textbf{意外删除}：全量仅 \textcolor{evidence}{2} 行删除，无结构性删除。
  \item \textbf{重命名}：\texttt{git diff -M} 未将任何「删+增」合并为重命名，以新增为主。
\end{itemize}

\section{结论与后续}

\begin{conclbox}
从 init 到 \texttt{stab27}，PyQCU 完成从空仓库到完整 Lattice QCD 库的建设：核心增量在 \texttt{cpp/}（CUDA 后端）与 \texttt{pyqcu/}（Python 实现 + Cython 桥 + 多线程多卡 MG），主导文件直接对应 Wilson/Clover dslash、多重网格与测试体系。
\end{conclbox}

\begin{refbox}
\texttt{refer/}（外部库镜像，1559 文件）与 \texttt{examples/*.txt}、\texttt{logs/} 属参考/产物，不应计入「库自有代码」规模评估。
\end{refbox}

\begin{warnbox}
后续建议：①将 \texttt{logs/}、\texttt{examples/pyqcu/log-*.txt} 等运行产物移出版本跟踪（\texttt{git rm --cached}）；②\texttt{refer/git-rep/} 体积庞大，建议以 submodule 或外部引用替代整库镜像；③聚焦「库代码」可用 \texttt{git diff 0204fa5 stab27 -- pyqcu/ cpp/} 排除噪声。
\end{warnbox}

\end{document}
